\documentclass[10pt]{article}
\usepackage{icmlko}

\icmlmeta{IP 파운데이션 월드모델}{162}{잣대를 다섯 재고 나서 칸이 모자란 걸 알았다}{손잡이 충실도를 재는 눈금 다섯을 후보로 놓고 폭·안정·판과의 반상관으로 골랐다. 제일 나은 것을 골랐는데, 그 폭이 이항 표준오차와 같았다}

\begin{document}

\twocolumn[
\icmlheader

\begin{abstract}
노트 161이 ``크기 가중이 다섯 정식화를 91$\sim$93\%에 몰아넣어 변별이
안 된다''를 남겼다. 눈금을 다섯 만들어 재 봤다 --- 단순 · 크기가중 ·
순위가중 · \textbf{상위절반}(효과가 큰 절반의 칸에서만 부호를 센다) ·
분해칸비율. 고르는 기준은 점수가 아니라 셋이다: \textbf{폭}(정식화 다섯
사이가 얼마나 벌어지나) · \textbf{안정}(씨앗을 바꿔도 순위가 같나) ·
\textbf{판과의 반상관}(노트 156이 세운 현상을 잡아내나). \textbf{상위절반이
셋 다 제일 낫다} --- 폭 0.056(크기가중의 2.5배) · 씨앗 순위 일치 $+$1.000 ·
판과 $-$0.866. 그 눈금으로 보면 그림이 아주 선명하다. \textbf{효과가 큰
열여덟 칸에서 선형 셋(F6 능형 · F9 순위우도 · F10 도메인수축)은
\emph{18/18} 로 다 맞히고 부스팅 둘(F18 배깅 · F8 부스팅)은 17/18 이다.}
노트 156부터 161까지 다섯 노트가 말해 온 ``판을 잘 맞힐수록 손잡이를 못
돌린다''가 여기서 제일 또렷하다. \textbf{그런데 여기서 멈추면 안 된다.}
18칸에서 $p{\approx}0.95$ 의 이항 표준오차가 0.051이고 $2\,\mathrm{se}$ 가
0.103이다. \textbf{관측된 폭 0.056은 1 se 다} --- 즉 18/18과 17/18의 차이는
\emph{한 칸}이고, 어떤 가중을 쓰든 이 표본으로는 못 가른다. 0.056을
$2\,\mathrm{se}$ 로 가르려면 큰 칸이 \textbf{60개} 필요한데 지금 18개다.
\textbf{병목은 잣대가 아니라 칸 수였다.} 칸은 도메인 $\times$ 축이고 지금
$9 \times 5$ 에서 관측된 것이 36, 그 절반이 18이다. 잣대는 상위절반으로
채택해 실행마다 적되, \textbf{정식화 사이 순위는 이 값으로 매기지 않는다.}
\end{abstract}
]

\begin{figure*}[t]
\centering
\includegraphics[width=\textwidth]{figs/ruler.pdf}
\caption{왼쪽 --- 상위절반이 셋 다 제일 낫다. 가운데 --- 선형 셋이 18/18.
오른쪽 --- 잣대가 아니라 칸이 모자라다.}
\end{figure*}

\section{눈금 다섯을 후보로 놓았다}

\begin{table}[t]
\centering\tsize
\caption{정식화 다섯에 대 봤다. 씨앗 둘로 안정도 잰다.}
\begin{tabular}{lrrr}
\toprule
& 폭 & 판과 상관 & 씨앗 순위 일치 \\
\midrule
단순 & \textbf{.139} & $-$.564 & $+$.973 \\
크기가중 & .022 & $-$.600 & $+$1.000 \\
순위가중 & .039 & $-$.500 & $+$.975 \\
\textbf{상위절반} & .056 & \textbf{$-$.866} & \textbf{$+$1.000} \\
분해칸비율 & .028 & \textbf{$-$.866} & $+$.968 \\
\bottomrule
\end{tabular}
\end{table}

\claimx{\textbf{상위절반이 셋 다 낫다.} 단순은 폭이 제일 크지만(0.139)
노트 159가 보인 대로 0.002짜리 칸을 0.06짜리와 같이 센다 --- 그래서 판과의
상관이 $-$0.564로 약하다. 크기가중은 폭이 0.022로 뭉갠다. 상위절반은 폭
0.056에 판과 $-$0.866, 씨앗 순위 완전 일치다.}{``분해칸비율''(효과가 0과
갈리는 칸의 비율)도 판과 $-$0.866인데 폭이 0.028이고 씨앗 안정이 조금
낮다. 게다가 이 판에서는 36칸 중 35칸이 이미 분해돼 변별이 거의 없다.}

\section{그 눈금으로 보면}

\begin{table}[t]
\centering\tsize
\caption{효과가 큰 열여덟 칸에서 사전 부호 일치.}
\begin{tabular}{lrrl}
\toprule
& 판 $\rho$ & 상위절반 & \\
\midrule
F18 배깅 & \textbf{$+$.4491} & 17/18 & .944 \\
F8 부스팅 & $+$.4341 & 17/18 & .944 \\
F10 도메인수축 & $+$.3735 & \textbf{18/18} & \textbf{1.000} \\
F6 능형 & $+$.3671 & \textbf{18/18} & \textbf{1.000} \\
F9 순위우도 & $+$.3519 & \textbf{18/18} & \textbf{1.000} \\
\bottomrule
\end{tabular}
\end{table}

\claimx{\textbf{선형 셋이 다 맞히고 부스팅 둘이 하나씩 틀린다.} 노트
156부터 다섯 노트가 말해 온 ``판을 잘 맞힐수록 손잡이를 못 돌린다''가
여기서 제일 또렷하다 --- 판 1 · 2위가 유일하게 틀린 둘이다.}{효과가 작은
나머지 열여덟 칸에서는 순서가 다르다(단순 기준 F6 · F9 89\%, F18 86\%,
F8 83\%, F10 75\%). \textbf{큰 칸과 작은 칸이 다른 이야기를 한다.}}

\section{그런데 폭이 표준오차다}

\begin{table}[t]
\centering\tsize
\caption{$p{\approx}0.95$ 에서 이항 표준오차.}
\begin{tabular}{lrr}
\toprule
큰 칸 & se & $2\,\mathrm{se}$ \\
\midrule
\textbf{18 (지금)} & \textbf{.051} & \textbf{.103} \\
36 & .036 & .073 \\
\textbf{60} & .028 & \textbf{.056} \\
100 & .022 & .044 \\
\bottomrule
\end{tabular}
\end{table}

\claimx{\textbf{관측된 폭 0.056은 1 se 다.} 18/18과 17/18의 차이는 한
칸이고, 어떤 가중을 써도 이 표본으로는 못 가른다. \textbf{병목은 잣대가
아니라 칸 수였다.}}{다섯 눈금을 다 재고 나서야 알았다 --- 재기 전에는
``크기 가중이 뭉갠다''가 눈금 탓인지 표본 탓인지 갈리지 않았다. \textbf{눈금
비교가 헛일은 아니었지만 답은 눈금 밖에 있었다.}}

\claimx{\textbf{60칸이 필요하다.} 지금 도메인 아홉 $\times$ 축 다섯에서
관측되는 것이 36칸이고 큰 절반이 18이다. 60까지 가려면 도메인이나 축을
늘려야 한다 --- 팝업 전용 다섯 축(체험 밀도 · 포토존 · 계절 적합 · 협업
강도 · IP 인지)을 넣으면 5칸이 는다. 부족하다.}{도메인을 늘리는 쪽이
크다. 도메인 하나가 최대 5칸이니 넷을 더 붙이면 56칸이다.}

\section{그래서 어떻게 쓰나}

\claimx{\textbf{상위절반을 실행마다 적되 순위는 이 값으로 안 매긴다.} 노트
144의 귀착, 146의 씨앗, 147의 짝 분해능과 같은 자리 --- 기록이지 판정이
아니다.}{``선형이 손잡이로 낫다''는 방향은 여섯 노트에 걸쳐 다섯 가지
다른 잣대에서 같은 쪽을 가리킨다. \textbf{한 잣대로는 못 가르지만 여섯
번 같은 방향이 나온 것은 다른 종류의 증거다.}}

\begin{table}[t]
\centering\tsize
\caption{남은 것.}
\begin{tabular}{ll}
\toprule
\textbf{칸} & 60칸이 필요하다 --- 도메인 넷을 더 붙이면 56 \\
큰/작은 칸 & 순서가 다르다 --- 무엇이 큰 칸을 만드나 \\
팝업 전용 축 & 다섯 더 있는데 공유가 아니라 안 쓰고 있다 \\
\bottomrule
\end{tabular}
\end{table}

둘째 줄이 곧바로 잴 수 있다. 큰 칸에서는 선형 셋이 완승이고 작은 칸에서는
순서가 섞인다. \textbf{무엇이 칸을 크게 만드는가} --- 그 도메인의 표본 수인지,
그 축의 분산인지, 라벨과의 상관인지. 셋 다 잴 수 있고 답에 따라 60칸을
어디서 구할지가 갈린다.

\end{document}
